\documentclass{exam-zh}
\usepackage{edu-practice}

\examsetup{
  question/show-answer = false,
  fillin/show-answer = false,
  solution/show-solution = hide,
}

% ---------- 教师版解答样式 ----------
\newtcolorbox{solutionblock}{
  enhanced, breakable,
  colback=edu-blue-bg,
  colframe=edu-blue!25,
  boxrule=0pt,
  borderline west={2pt}{0pt}{edu-blue!50},
  arc=0pt,
  left=3mm, right=3mm, top=2mm, bottom=2mm,
  before skip=2mm,
  after skip=3mm,
  fontupper=\small,
}

% 解答步骤编号
\newcounter{solstep}
\newcommand{\solstep}[2]{%
  \stepcounter{solstep}%
  \par\medskip
  \noindent\hangindent=6mm
  {\color{edu-blue}\bfseries\textcircled{\scriptsize\thesolstep}\enspace #1}\enspace #2\par
}

\begin{document}

\section*{立体几何定理库默写 · 全覆盖竖式版}

\section*{一、点线面基本事实与分类}


\begin{question}[points=5]
  \textbf{符号约定}\par
\[
\left.
\begin{aligned}
A&\ \underline{\hspace{9mm}}\ l\\
l&\ \underline{\hspace{9mm}}\ \alpha\\
\alpha&\ \underline{\hspace{9mm}}\ \beta=m
\end{aligned}
\right\}
\Longrightarrow
\begin{aligned}
&\text{点在线上；}\\
&\text{线在面内；}\\
&\text{两面交于直线 }m
\end{aligned}
\]


\end{question}



\begin{question}[points=5]
  \textbf{位置关系分类}\par
\[
\left.
\begin{aligned}
\text{线线}&:\ \text{相交、平行、}\underline{\hspace{12mm}}\\
\text{线面}&:\ \text{在面内、相交、}\underline{\hspace{12mm}}\\
\text{面面}&:\ \text{平行、}\underline{\hspace{12mm}}
\end{aligned}
\right\}
\Longrightarrow \text{先判对象，再选定理}
\]


\end{question}



\begin{question}[points=5]
  \textbf{基本事实1：三点定面}\par
\[
\left.
\begin{aligned}
A,B,C&\ \underline{\hspace{18mm}}
\end{aligned}
\right\}
\Longrightarrow
\text{过 }A,B,C\text{ 有且只有一个平面}
\]


\end{question}



\begin{question}[points=5]
  \textbf{基本事实2：点线面从属}\par
\[
\left.
\begin{aligned}
A,B&\in l\\
A,B&\in\alpha\\
A&\ne B
\end{aligned}
\right\}
\Longrightarrow
\underline{\hspace{24mm}}
\]


\end{question}



\begin{question}[points=5]
  \textbf{基本事实3：公共点推出交线}\par
\[
\left.
\begin{aligned}
P&\in\alpha\\
P&\in\beta\\
\alpha&\ne\beta
\end{aligned}
\right\}
\Longrightarrow
\alpha\cap\beta=\underline{\hspace{12mm}},\quad P\in\underline{\hspace{10mm}}
\]


\end{question}



\begin{question}[points=5]
  \textbf{推论1：直线与线外一点定面}\par
\[
\left.
\begin{aligned}
a&\text{ 是直线}\\
P&\notin a
\end{aligned}
\right\}
\Longrightarrow
\underline{\hspace{42mm}}
\]


\end{question}



\begin{question}[points=5]
  \textbf{推论2：相交直线定面}\par
\[
\left.
\begin{aligned}
a\cap b&=P
\end{aligned}
\right\}
\Longrightarrow
\underline{\hspace{38mm}}
\]


\end{question}



\begin{question}[points=5]
  \textbf{推论3：平行直线定面}\par
\[
\left.
\begin{aligned}
a&\parallel b\\
a&\ne b
\end{aligned}
\right\}
\Longrightarrow
\underline{\hspace{38mm}}
\]


\end{question}


\section*{二、平行关系定理}


\begin{question}[points=5]
  \textbf{平行传递性}\par
\[
\left.
\begin{aligned}
a&\parallel b\\
b&\parallel c
\end{aligned}
\right\}
\Longrightarrow
\underline{\hspace{20mm}}
\]


\end{question}



\begin{question}[points=5]
  \textbf{等角定理}\par
\[
\left.
\begin{aligned}
a&\parallel a'\\
b&\parallel b'
\end{aligned}
\right\}
\Longrightarrow
\angle(a,b)\text{ 与 }\angle(a',b')\ \underline{\hspace{26mm}}
\]


\end{question}



\begin{question}[points=5]
  \textbf{线面平行判定}\par
\[
\left.
\begin{aligned}
a&\parallel b\\
b&\subset\alpha\\
a&\not\subset\alpha
\end{aligned}
\right\}
\Longrightarrow
\underline{\hspace{22mm}}
\]


\end{question}



\begin{question}[points=5]
  \textbf{线面平行性质}\par
\[
\left.
\begin{aligned}
a&\parallel\alpha\\
a&\subset\beta\\
\alpha\cap\beta&=b
\end{aligned}
\right\}
\Longrightarrow
\underline{\hspace{22mm}}
\]


\end{question}



\begin{question}[points=5]
  \textbf{面面平行判定}\par
\[
\left.
\begin{aligned}
a,b&\subset\alpha\\
a\cap b&=P\\
a&\parallel\beta\\
b&\parallel\beta
\end{aligned}
\right\}
\Longrightarrow
\underline{\hspace{22mm}}
\]


\end{question}



\begin{question}[points=5]
  \textbf{面面平行性质}\par
\[
\left.
\begin{aligned}
\alpha&\parallel\beta\\
\gamma\cap\alpha&=a\\
\gamma\cap\beta&=b
\end{aligned}
\right\}
\Longrightarrow
\underline{\hspace{22mm}}
\]


\end{question}



\begin{question}[points=5]
  \textbf{面面平行的常用性质}\par
\[
\left.
\begin{aligned}
\alpha&\parallel\beta\\
a&\subset\alpha
\end{aligned}
\right\}
\Longrightarrow
\underline{\hspace{22mm}}
\]


\end{question}


\section*{三、垂直关系定理}


\begin{question}[points=5]
  \textbf{异面直线垂直定义}\par
\[
\left.
\begin{aligned}
a'&\parallel a\\
a'\cap b&=O\\
\angle(a',b)&=90^\circ
\end{aligned}
\right\}
\Longrightarrow
\underline{\hspace{20mm}}
\]


\end{question}



\begin{question}[points=5]
  \textbf{线面垂直定义推论}\par
\[
\left.
\begin{aligned}
a&\perp\alpha\\
b&\subset\alpha
\end{aligned}
\right\}
\Longrightarrow
\underline{\hspace{20mm}}
\]


\end{question}



\begin{question}[points=5]
  \textbf{线面垂直判定}\par
\[
\left.
\begin{aligned}
l&\perp m\\
l&\perp n\\
m,n&\subset\alpha\\
m\cap n&=P
\end{aligned}
\right\}
\Longrightarrow
\underline{\hspace{20mm}}
\]


\end{question}



\begin{question}[points=5]
  \textbf{线面垂直性质}\par
\[
\left.
\begin{aligned}
a&\perp\alpha\\
b&\perp\alpha
\end{aligned}
\right\}
\Longrightarrow
\underline{\hspace{20mm}}
\]


\end{question}



\begin{question}[points=5]
  \textbf{面面垂直判定}\par
\[
\left.
\begin{aligned}
l&\perp\alpha\\
l&\subset\beta
\end{aligned}
\right\}
\Longrightarrow
\underline{\hspace{20mm}}
\]


\end{question}



\begin{question}[points=5]
  \textbf{面面垂直性质}\par
\[
\left.
\begin{aligned}
\alpha&\perp\beta\\
\alpha\cap\beta&=l\\
a&\subset\alpha\\
a&\perp l
\end{aligned}
\right\}
\Longrightarrow
\underline{\hspace{20mm}}
\]


\end{question}



\begin{question}[points=5]
  \textbf{三垂线定理}\par
\[
\left.
\begin{aligned}
PA&\perp\alpha\\
AB&\subset\alpha\\
l&\subset\alpha\\
AB&\perp l
\end{aligned}
\right\}
\Longrightarrow
\underline{\hspace{20mm}}
\]


\end{question}



\begin{question}[points=5]
  \textbf{三垂线逆定理}\par
\[
\left.
\begin{aligned}
PA&\perp\alpha\\
AB&\subset\alpha\\
l&\subset\alpha\\
PB&\perp l
\end{aligned}
\right\}
\Longrightarrow
\underline{\hspace{20mm}}
\]


\end{question}



\begin{question}[points=5]
  \textbf{最小角定理}\par
\[
\left.
\begin{aligned}
l\cap\alpha&=A\\
AH&\text{ 是 }l\text{ 在 }\alpha\text{ 内的射影}\\
b&\subset\alpha,\ A\in b
\end{aligned}
\right\}
\Longrightarrow
\angle(l,AH)\ \underline{\hspace{16mm}}\ \angle(l,b)
\]


\end{question}


\section*{四、距离与角的化归规则}


\begin{question}[points=5]
  \textbf{点到直线距离}\par
\[
\left.
\begin{aligned}
PH&\perp l\\
H&\in l
\end{aligned}
\right\}
\Longrightarrow
d(P,l)=\underline{\hspace{16mm}}
\]


\end{question}



\begin{question}[points=5]
  \textbf{点到平面距离}\par
\[
\left.
\begin{aligned}
PH&\perp\alpha\\
H&\in\alpha
\end{aligned}
\right\}
\Longrightarrow
d(P,\alpha)=\underline{\hspace{16mm}}
\]


\end{question}



\begin{question}[points=5]
  \textbf{线面距转点面距}\par
\[
\left.
\begin{aligned}
a&\parallel\alpha\\
P&\in a
\end{aligned}
\right\}
\Longrightarrow
d(a,\alpha)=\underline{\hspace{20mm}}
\]


\end{question}



\begin{question}[points=5]
  \textbf{面面距转点面距}\par
\[
\left.
\begin{aligned}
\alpha&\parallel\beta\\
P&\in\alpha
\end{aligned}
\right\}
\Longrightarrow
d(\alpha,\beta)=\underline{\hspace{20mm}}
\]


\end{question}



\begin{question}[points=5]
  \textbf{异面直线距离}\par
\[
\left.
\begin{aligned}
M&\in a,\ N\in b\\
MN&\perp a\\
MN&\perp b
\end{aligned}
\right\}
\Longrightarrow
d(a,b)=\underline{\hspace{16mm}}
\]


\end{question}



\begin{question}[points=5]
  \textbf{等体积法求点面距}\par
\[
\left.
\begin{aligned}
V_{P-ABC}&=\frac13 S_{\triangle ABC}\cdot h
\end{aligned}
\right\}
\Longrightarrow
h=\underline{\hspace{28mm}}
\]


\end{question}



\begin{question}[points=5]
  \textbf{异面直线所成角}\par
\[
\left.
\begin{aligned}
a'&\parallel a\\
b'&\parallel b\\
a'\cap b'&=O
\end{aligned}
\right\}
\Longrightarrow
\angle(a,b)=\underline{\hspace{24mm}}
\]


\end{question}



\begin{question}[points=5]
  \textbf{线面角定义}\par
\[
\left.
\begin{aligned}
l\cap\alpha&=A\\
P&\in l\\
P\text{ 在 }\alpha\text{ 上的射影为 }H
\end{aligned}
\right\}
\Longrightarrow
\angle(l,\alpha)=\underline{\hspace{18mm}}
\]


\end{question}



\begin{question}[points=5]
  \textbf{点面距法求线面角}\par
\[
\left.
\begin{aligned}
\theta&=\angle(l,\alpha)\\
P&\in l,\ A=l\cap\alpha\\
d(P,\alpha)&=d\\
PA&=s
\end{aligned}
\right\}
\Longrightarrow
\sin\theta=\underline{\hspace{18mm}}
\]


\end{question}



\begin{question}[points=5]
  \textbf{二面角平面角定义}\par
\[
\left.
\begin{aligned}
\alpha\cap\beta&=l\\
O&\in l\\
OA&\subset\alpha,\ OB\subset\beta\\
OA&\perp l,\ OB\perp l
\end{aligned}
\right\}
\Longrightarrow
\underline{\hspace{16mm}}\text{ 是二面角 }\alpha-l-\beta\text{ 的平面角}
\]


\end{question}



\begin{question}[points=5]
  \textbf{二面角面积射影法}\par
\[
\left.
\begin{aligned}
S_{\text{射影}}&=S_{\text{原}}\cos\theta
\end{aligned}
\right\}
\Longrightarrow
\cos\theta=\underline{\hspace{28mm}}
\]


\end{question}



\end{document}
